\BioseqWalkRow{sql}{appendices/02_sql_walkthrough/sql/01_create.sql}{%
\textbf{1. Create a sequence table.}
The four ordinary columns hold patient metadata. The sequence column uses
the DNA4 alphabet, which accepts A, C, G and T. The number 64 is the packed
storage-unit width, not a sequence-length limit.

\texttt{USING bioseq} selects BioSeqDB storage and search planning.
Figure~\ref{fig:storage-heap-to-stripe} shows how its ordinary columns are
stored. Section~\ref{sec:sequence-bit-layout} explains sequence packing.
The first two lines tell \texttt{psql} to stop on errors and avoid paging.}

\BioseqWalkRow{sql}{appendices/02_sql_walkthrough/sql/02_insert.sql}{%
\textbf{2. Insert three related sequences.}
The common expression makes a reproducible 2,048-base string by translating
hexadecimal digits to DNA letters. This is synthetic data, not an MD5-based
biological model.

Anna receives two copies of that string, Ben a 512-base piece, and Cara one
copy. Their lengths are 4,096, 512 and 2,048 bases. Ben's sequence occurs
twice in Anna's sequence, which will make multiple alignment rows easy to see.
PostgreSQL's text \texttt{substring} starts at 1, so position 257 here
corresponds to BioSeqDB position 256.}

\BioseqWalkRow{sql}{appendices/02_sql_walkthrough/sql/03_select.sql}{%
\textbf{3. Read ordinary columns.}
This returns all three patients. No sequence symbols are needed, because
the projection contains only metadata.
Without \texttt{ORDER BY}, SQL does not promise a row order.}

\BioseqWalkRow{sql}{appendices/02_sql_walkthrough/sql/04_filter.sql}{%
\textbf{4. Combine metadata conditions and a sequence preview.}
This returns Ben and Anna, in that order. The age interval, date and name
conditions must all hold.

\texttt{len} reports sequence symbols. The slice \texttt{[:16]} returns
the first sixteen symbols and the text cast makes them readable.
BioSeqDB slices use zero-based, half-open bounds, as shown in
Figure~\ref{fig:sequence-slice-coordinates}.}

\BioseqWalkRow{sql}{appendices/02_sql_walkthrough/sql/05_explain_read.sql}{%
\textbf{5. Explain the unfiltered read.}
\texttt{ANALYZE} executes statement 3 and measures it.
\texttt{COSTS OFF} hides planner estimates to keep this first plan small.
It does not disable execution measurements.}

\BioseqWalkRow{plan}{appendices/02_sql_walkthrough/plans/05_explain_read.plan}{%
\texttt{BioseqScan} reads the table. \texttt{Projected Columns} lists
exactly the four requested columns. \texttt{actual time} gives milliseconds
to the first and last output row, averaged per loop. Here there are three
rows and one loop. Planning and execution times describe the whole statement.

The cache stores decompressed column chunks, as described in
Table~\ref{tab:storage-cache-scopes}. Hits and misses count buffer lookups.
Writes count buffers added or replaced, evictions count removed entries,
and entries count current buffers. Maximum and ending sizes are peak and
current payload bytes, not the configured cache capacity. These counters
belong to the backend's cache state, not to a cache of SQL results.

There is no chunk-skipping field here because there is no filter. The
filtered search below shows that field.}

\BioseqWalkRow{sql}{appendices/02_sql_walkthrough/sql/06_search.sql}{%
\textbf{6. Search with a simple metadata filter.}
\texttt{\textbackslash gset} saves Ben's DNA in the \texttt{psql} variable
\texttt{probe}. The spelling \texttt{:'probe'} safely inserts it as an SQL
string literal. The following constant-query searches reuse this
512-base query.

The stored target is the first \texttt{mmsearch} argument. This returns
two matches for Anna and one for Ben. Each has score 1,024 and CIGAR
\texttt{512M}.

Requesting \texttt{cigar} enables base alignment. In this alpha version,
\texttt{score} alone does not enable it and can return zero.
Section~\ref{sec:final-alignment-role} explains that stage.
For the misleading allocation error, see
Warning~\ref{warn:walkthrough-alpha}.}

\BioseqWalkRow{sql}{appendices/02_sql_walkthrough/sql/07_build.sql}{%
\textbf{7. Build the MPF index explicitly.}
This processes three sequence values and returns 3. It extracts minimizers
using the current algorithm and $k,w$, groups their positions, and publishes
the relation data and its reader files. See
Algorithm~\ref{alg:index-image-lifecycle}.

The default grouped format requires the two asynchronous-persistence
settings to be off. They already are off, so this example needs no
preparatory \texttt{SET} commands. Update policies appear in
Table~\ref{tab:index-update-policy}.}

\BioseqWalkRow{sql}{appendices/02_sql_walkthrough/sql/08_search_again.sql}{%
\textbf{8. Repeat the search after construction.}
The same three matches remain. A successful index build does not imply
that every query reads its file. This ordinary constant-query scan
recomputes target minimizers in the current default configuration.

The two-table join later uses the required persisted-reader path.
Figure~\ref{fig:index-join-read-paths} shows the distinction.}

\BioseqWalkRow{sql}{appendices/02_sql_walkthrough/sql/09_options.sql}{%
\textbf{Choose algorithms for one call.}
These are the same defaults written explicitly to demonstrate
\texttt{aligner\_options}, so the result still has three matches.
The codes choose bounded ntHash SIMD minimizers (4), SIMD minimap2
chaining (1), and BlockAligner (1). Each minimizer window contains ten
consecutive 15-mers.

Omitted options inherit session settings. Appendix~\ref{app:configuration-reference}
lists the numeric codes beside their named GUC values.
Section~\ref{sec:arch-request-settings} explains precedence.}

\BioseqWalkRow{sql}{appendices/02_sql_walkthrough/sql/10_explain_search.sql}{%
\textbf{9. Explain exactly that search.}
The following consecutive blocks contain its complete text plan.
Keeping the CIGAR projection matters, because removing it would also
remove the request for base alignment.}

\BioseqWalkRow[firstline=1,lastline=8]{plan}{appendices/02_sql_walkthrough/plans/10_explain_search.plan}{%
\texttt{Result} projects the output from \texttt{BioseqScan}. Three
alignment rows come from two qualifying patients.

The scan reads the identifier, age and sequence. A chunk-group filter
compares the age condition with stored minima and maxima. Zero groups
were skipped here. The vectorized filter tests individual ages in the
remaining batches. Listing~\ref{lst:storage-minmax-pruning} explains this
difference.

\texttt{allow} permits ordinary index reuse or recomputation.
The planned share is 100\% recomputed. The actual counters below confirm
it. See Table~\ref{tab:index-read-requirement}.}

\BioseqWalkRow[firstline=9,lastline=18,gobble=8]{plan}{appendices/02_sql_walkthrough/plans/10_explain_search.plan}{%
The scheduler line is a snapshot of the local task scheduler.
It reports workers, processed/scheduled batches, waiting work and
processed/waiting queries. The local supervisor leaves the query-rate
field at zero. The last fields count pending persistence jobs and reserved
index bytes. These are not timings for individual alignment stages.
Section~\ref{sec:arch-work-distribution} explains its queues.

Highway selects a CPU instruction set. Runnable targets are the intersection
of compiled and CPU-supported targets. Planned and executed architectures
agree here. Your CPU may show different names.
\texttt{QuerySketchStage} extracts the query minimizers, and
\texttt{Highway Function} names its selected kernel.
See Section~\ref{sec:arch-hardware-selection}.}

\BioseqWalkRow[firstline=19,lastline=27,gobble=12]{plan}{appendices/02_sql_walkthrough/plans/10_explain_search.plan}{%
These are internal pipeline descriptions.
The zero \texttt{cep.window\_len}, \texttt{cep.kmer\_len} and
\texttt{query\_len} values are placeholders in that description.
They do not mean that this query has length zero or uses $k=w=0$.
\texttt{sequence\_len\_log2} is a compiled size category.

\texttt{specialised: fallback} means a general parameter implementation,
not a failed GPU request. \texttt{max\_occ} limits accepted occurrences
of one minimizer hash. \texttt{reserved} is a planned anchor-buffer capacity
in elements, not a byte count or the number of matches.}

\BioseqWalkRow[firstline=28,lastline=41,gobble=12]{plan}{appendices/02_sql_walkthrough/plans/10_explain_search.plan}{%
Chaining joins seed anchors into candidate regions.
\texttt{min\_chain\_score} and \texttt{min\_seed\_count} reject weak
chains. \texttt{max\_chain\_skip} and \texttt{max\_chain\_iter} limit
predecessor exploration. Bandwidth and maximum gap constrain how far
anchors may separate.

This is the bounded SIMD approximation described in
Section~\ref{sec:chaining-bioseqdb}. Its flags describe the vectorized implementation,
whether a zero skip penalty is required, and the general specialization.
The next stage makes regions from those chains.}

\BioseqWalkRow[firstline=42,lastline=49,gobble=12]{plan}{appendices/02_sql_walkthrough/plans/10_explain_search.plan}{%
Regions are classified as primary or secondary by their overlap.
\texttt{mask\_level} is the overlap threshold and \texttt{mask\_len}
limits uncovered query length. \texttt{sub\_diff} is the close-score
margin used when counting competitors.

Secondary filtering uses a score ratio, an absolute score margin
(\texttt{min\_diff}), and \texttt{best\_n}. The last value limits retained
secondary regions, not SQL result rows or primary matches.
See Section~\ref{sec:alignment-postprocessing}.}

\BioseqWalkRow[firstline=50,lastline=60,gobble=12]{plan}{appendices/02_sql_walkthrough/plans/10_explain_search.plan}{%
BlockAligner computes the base-level alignment and CIGAR.
Matching bases add 2 to the score, substitutions cost 4, and the displayed
gap-opening and extension penalties define its single affine-gap scheme.
Its adaptive block size ranges from 32 to 256.

The vectorized and single-affine flags identify this implementation.
The Highway line identifies its CPU kernel.
See Section~\ref{sec:final-alignment-cpu-contracts}.}

\BioseqWalkRow[firstline=61,lastline=69,gobble=12]{plan}{appendices/02_sql_walkthrough/plans/10_explain_search.plan}{%
Primary/secondary classification and filtering run again after base
alignment has refined the regions. Their parameters have the same meanings
as above.

\texttt{Submit results} reports two processed targets. Both used recomputed
minimizers and neither used a persisted reader. This confirms why the
preceding index build did not change this scan's chosen read path.}

\BioseqWalkRow[firstline=70,lastline=72,gobble=12]{plan}{appendices/02_sql_walkthrough/plans/10_explain_search.plan}{%
These diagnostics try to certify which implementations completed the work.
\texttt{attempted}, \texttt{fully\_attested} and \texttt{unattested}
count attempted and certified target/query jobs. The first uncertified
target and query ordinal are printed below them.

\texttt{mixed\_invalid} with two unattested targets means the current
instrumentation did not certify a complete backend combination.
It is not evidence of GPU execution or a scalar fallback.
The flags report mixed environments, mixed persisted identities, conflicting
identity evidence and saturated counters. All are false here.

Requested and actual index sources differ. \texttt{missing\_persisted}
means no suitable persisted reader was used for these targets, even though
physical MPF data exists. The permissive policy allows recomputation.}

\BioseqWalkRow[firstline=73,lastline=77,gobble=12]{plan}{appendices/02_sql_walkthrough/plans/10_explain_search.plan}{%
GPU counters cover seed lookup, chaining, backtracking, coordinate voting
and extension. Declines and fallbacks describe execution-route decisions.
All are zero in this CPU example. Three counters are also printed under
their longer names.

Host statistics include query-sketch and target-build nanoseconds, hash
readbacks, symbol-packing calls, alignment classification and CIGAR
construction. Zeros in this incompletely attested execution do not mean
that these operations took no time or that no CIGAR was produced.}

\BioseqWalkRow[firstline=78,lastline=87,gobble=8]{plan}{appendices/02_sql_walkthrough/plans/10_explain_search.plan}{%
The cache fields have the same meanings as in the first plan. This search
needs sequence data as well as metadata, so the buffers and their sizes
differ. Cache entries count column chunks, not matches.

The final two lines give statement planning and execution time in
milliseconds. They include work outside the alignment stages.}

\BioseqWalkRow{sql}{appendices/02_sql_walkthrough/sql/11_self.sql}{%
\textbf{10. Search each sequence within itself.}
Each of the three rows produces a hit. With CIGAR requested, these exact
self-alignments cover their complete sequences and score twice their length.

This is a same-row search. It performs alignment rather than replacing
\texttt{mmsearch(seq,seq)} with a constant Boolean value.}

\BioseqWalkRow{sql}{appendices/02_sql_walkthrough/sql/12_self_join.sql}{%
\textbf{Compare rows using a self-join.}
The two aliases independently choose target and query rows. The result has
13 alignments, including the three equal-ID pairs.
Section~\ref{sec:mmsearch-join-ownership} explains these roles.

For target 1 and query 2, score 1,024 occurs at both $[256,768)$ and
$[2304,2816)$. These are different occurrences, even though both have
\texttt{512M}. Keep positions as well as scores.
All helper arguments name the target \texttt{t.seq}, including the query
coordinate helpers. Figure~\ref{fig:scan-occurrences} explains why a stored
row can produce several output rows.}

\BioseqWalkRow[firstline=1,lastline=14]{sql}{appendices/02_sql_walkthrough/sql/13_fasta.sql}{%
\textbf{11. Load a FASTA file into another table.}
The first statement writes our three records as alternating header and
sequence lines. This makes the exercise independent of external downloads.
The server must be able to write and then read the named path.

\texttt{load} returns 3. With one text column, one sequence column and no
explicit mappings, its default \texttt{auto} mode keeps each full header
without the leading \texttt{>}. It does not create an identity column or
invent numeric IDs.}

\BioseqWalkRow[firstline=15,lastline=16]{sql}{appendices/02_sql_walkthrough/sql/13_fasta.sql}{%
\textbf{Header parsing options.}
This call loads the file once. Running it again would append duplicate
records.

\texttt{header\_format} accepts \texttt{auto}, \texttt{plain},
\texttt{uni\_prot}, \texttt{uni\_ref}, \texttt{uni\_parc},
\texttt{uni\_mes}, \texttt{spd}, \texttt{ncbi}, \texttt{ncbi\_gi},
\texttt{ncbi\_pdb}, \texttt{refseq}, \texttt{genbank}, \texttt{gisaid}
and \texttt{imgt}. Auto can detect a dialect when field parsing is needed.
Plain keeps an unsplit identifier.

The fourth argument maps parsed fields to table columns using
\texttt{bioseq.column\_mapping}, with destination first and source second.
For example, \texttt{('accession','accession')} copies that parsed field.
Section~\ref{sec:io-header-mapping} explains the mappings and loader limits.}

\BioseqWalkRow{sql}{appendices/02_sql_walkthrough/sql/14_loaded_search.sql}{%
\textbf{12. Inspect prefixes, then search a loaded slice.}
The first query returns three headers, lengths and sixteen-base previews.
The next saves a 256-base slice of the loaded Ben record as
\texttt{short\_probe}.

Searching it returns four alignments: two in Anna's sequence and one in
each of the other records. The search query is still a typed DNA4 value.
Taking its slice before the search avoids displaying a long literal.}

\BioseqWalkRow[firstline=1,lastline=15]{sql}{appendices/02_sql_walkthrough/sql/15_all_fields.sql}{%
\textbf{13. Return every alignment field.}
This statement continues in the next block. Scores and coordinates describe
the current match. \texttt{score} and \texttt{dp\_score} are aliases.
On this CPU path, \texttt{dp\_max\_score} currently repeats that score.
\texttt{chain\_score} describes the seed chain and the secondary score
describes a competitor.

Lengths count logical symbols. Both query and target coordinates use
zero-based, half-open intervals. \texttt{block\_length} is the alignment
block length.

\texttt{num\_matches} is a reported aligned-residue count, not a result-row
count. The current CPU code counts CIGAR \texttt{M} and \texttt{=} runs.
Since \texttt{M} can include substitutions, this field is not a reliable
count of identical bases for an imperfect alignment.}

\BioseqWalkRow[firstline=16,lastline=23]{sql}{appendices/02_sql_walkthrough/sql/15_all_fields.sql}{%
\texttt{edit} counts alignment edits. \texttt{mapq} estimates mapping
quality, while \texttt{divergence} is estimated from seeds. Neither is a
percentage identity. \texttt{seed\_count} counts chain anchors.

The Boolean fields report reverse orientation, primary classification and
whether CIGAR is present. CIGAR describes runs of alignment operations.
\texttt{M} consumes both sequences and can include a substitution.
\texttt{I} consumes the query, and
\texttt{D} consumes the target.
Section~\ref{sec:match-placeholder-injection} explains how these helpers
receive the current match's values.}

\BioseqWalkRow[firstline=24,lastline=28]{sql}{appendices/02_sql_walkthrough/sql/15_all_fields.sql}{%
This join requires the compatible persisted target index built in step 7.
Code 2 is a per-call read policy, not a GUC. It lets this supported
\texttt{BioseqBatchJoin} use the grouped MPF reader files.
A missing or incompatible reader is an error, as explained in
Section~\ref{sec:index-read-policy}.

The result has 13 alignments, ordered by descending score and then by
identifiers and positions. The complete statement includes all 21 helpers.}

\BioseqWalkRow{sql}{appendices/02_sql_walkthrough/sql/16_vector.sql}{%
\textbf{14. Aggregate metadata in batches.}
This returns two patients, minimum age 34, maximum age 52, sum 86 and
average 43. The dates span 1--2 September.

The eligible aggregate functions operate directly on ordinary integer and
date columns. This query does not need sequence symbols or alignments.}

\BioseqWalkRow{sql}{appendices/02_sql_walkthrough/sql/17_explain_vector.sql}{%
Explain the same aggregate. The excerpt below keeps only the node and
filter fields relevant to vectorization. The dots stand for the omitted
cache and timing lines, not additional SQL to execute.}

\BioseqWalkRow{plan}{appendices/02_sql_walkthrough/plans/vector_excerpt.plan}{%
\texttt{VectorAggNode} is the vector aggregate. Its child supplies column
batches and applies the vectorized age/date filter. Here the child reports
one batch slot, while the aggregate returns one summary row. Neither
\texttt{rows=1} means that only one patient was included.

This node currently has no separate per-lane or SIMD-speedup counter.
The analyzed node and its batch-reading child show the chosen execution
path. Grouping or aggregating alignment results can select an ordinary
aggregate instead.
Section~\ref{sec:filter-placement} distinguishes batch filtering from
other execution stages.}

\BioseqWalkRow{sql}{appendices/02_sql_walkthrough/sql/18_write.sql}{%
\textbf{15. Update, insert and rebuild.}
Anna becomes 35 and Daria receives Ben's sequence. Both writes return their
affected row. There are now four stored sequences.

With insert maintenance off, these writes do not build new minimizer
generations. Even a metadata update creates a replacement row version.
The two explicit calls prepare MPF and symbol-count indexes for all four
visible values. See Figure~\ref{fig:storage-update-snapshots} and
Section~\ref{sec:index-file-consistency}.}

\BioseqWalkRow{sql}{appendices/02_sql_walkthrough/sql/19_stats.sql}{%
Inspect the three registered index families. MPF and symbol counts have
stored data. MPS remains unbuilt because this exercise never requested it.

\texttt{storage\_size} measures the index relations and excludes separate
MPF reader files. Physical rows count internal index records, not patients.
The update age is the time since the recorded build.
\texttt{UPDATE\_ON\_INSERT} is the family's policy, which still needs the
separate maintenance switch. See Table~\ref{tab:index-policy-defaults}.}

\BioseqWalkRow{sql}{appendices/02_sql_walkthrough/sql/20_scalar.sql}{%
\textbf{Try scalar minimizers, scalar chaining and KSW2.}
The three algorithm codes are all 0. The query returns five matches,
each scoring 1,024 with \texttt{512M}.

Read requirement 1 forbids persisted target-index reads for this call.
That matters because the stored index was built with a different minimizer
algorithm. It remains stored for later compatible queries.
This local override does not change session defaults.
See Section~\ref{sec:index-compatibility}.}

\BioseqWalkRow{sql}{appendices/02_sql_walkthrough/sql/21_extras.sql}{%
\textbf{Combine sequence operations and an edit filter.}
This returns four alignment rows for Anna, Ben and Daria.
It counts A bases, joins two eight-base slices, reads one letter, and
reverses or reverse-complements a prefix. Reversing the reverse complement
gives the ordinary complement. There is no separate SQL complement helper.

The exact predicate \texttt{bioseq.edit(seq) < 5} permits an early edit
prefilter and still checks the retained alignment's edit count.
Use the target sequence column here, not its ID. Prefilter dispatch depends
on the supported candidate shape, as explained in
Section~\ref{sec:prefilter-role}.

The four matches have zero edits. The sequence expressions remain typed
until their final text casts, following Section~\ref{sec:sequence-sql-slicing}.}
